\appendix
\section{Algebraic identity sheet}
\label{app:identities}

This appendix collects the equations that a formal development can use as rewrite targets.

\subsection{Endpoint and center equations}

\begin{align*}
 t&=E-A=D-B,\\
 D&=B+t, & E&=A+t,\\
 p&=B+C, & q&=C+D=p+t,\\
 r&=E+F, & s&=F+A=r-t,\\
 u&=r-p=E+F-B-C.
\end{align*}

\subsection{Half-turn compositions}

For \(H_z(x)=z-x\),
\begin{align*}
 H_zH_w&=\tau_{z-w},\\
 H_rH_p&=\tau_u,\\
 H_rH_q&=\tau_{u-t},\\
 H_rH_r&=\mathrm{id},\\
 H_rH_s&=\tau_t.
\end{align*}

\subsection{Boundary equations}

\begin{align*}
 \tau_t(a)&=-d,\\
 H_p(b)&=-b, & H_r(b)&=-\tau_u(b),\\
 H_q(c)&=-c, & H_r(c)&=-\tau_{u-t}(c),\\
 H_r(e)&=-e,\\
 H_s(f)&=-f, & H_r(f)&=-\tau_t(f).
\end{align*}

\subsection{Shoelace equations}

For \(P:X\to Y\),
\begin{align*}
 \PathJ(\tau_wP)&=\PathJ(P)+\det(w,Y-X),\\
 \PathJ(H_zP)&=\PathJ(P)+\det(z,X-Y),\\
 \PathJ(-P)&=-\PathJ(P),\\
 H_{X+Y}P=-P&\Longrightarrow \PathJ(P)=\det(X,Y).
\end{align*}
Thus
\begin{align*}
 \PathJ(a)+\PathJ(d)&=-\det(t,B-A),\\
 \PathJ(b)+\PathJ(c)+\PathJ(e)+\PathJ(f)
 &=\det(B,C)+\det(C,D)+\det(E,F)+\det(F,A),\\
 \PathJ(\partial T)&=\det(t,u)=2\Area(T).
\end{align*}

\section{Primitive-recursive torus constructor}
\label{app:algorithm}

The following pseudocode spells out one total implementation.  Every loop has an explicit finite bound.

\begin{lstlisting}
BuildTorus(code c):
  if not DiskOK(c) or not ConwayOK(c):
      return EmptyCertificate

  t := E(c) - A(c)
  u := E(c) + F(c) - B(c) - C(c)
  delta := det(t,u)
  if delta <= 0: return EmptyCertificate

  triangles := DecodeCertifiedTriangles(c)
  turned := map (x |-> E(c)+F(c)-x) triangles
  pieces := []

  for triangle q in triangles ++ turned:
      q' := map LatticeCoordinates(t,u,delta) q
      I := all integer horizontal shifts whose box meets [0,1]^2
      J := all integer vertical shifts whose box meets [0,1]^2
      for i in I:
        for j in J:
          p := ClipConvex(q' + (i,j), UnitSquare)
          if TwiceArea(p) > 0: append p to pieces

  intersections := all pairwise segment intersections in pieces
  split every edge at its sorted intersection list
  graph := build directed half-edge graph with torus side tags
  faces := bounded face walk from every unused half-edge
  trianglesK := triangulate each rational face
  dualTree := least-search spanning tree of face adjacency
  counts := coverage count of each triangleK by original pieces

  return Certificate(pieces, graph, faces, trianglesK, dualTree, counts)
\end{lstlisting}

The verifier checks:
\begin{enumerate}[label=(\arabic*)]
\item all decoded rational denominators are positive;
\item all segment splits and cyclic face walks are exact;
\item every directed half-edge has the prescribed reverse, including side identifications;
\item every small triangle is positively oriented;
\item the dual tree visits every face exactly once;
\item the boundary coefficient of every edge is zero;
\item adjacent face counts agree; and
\item the common count is one by the determinant-area identity.
\end{enumerate}

A simpler certificate can omit \code{graph}, \code{faces}, and \code{counts}, retaining the clipped convex pieces and checking pairwise zero-area intersections plus total area one.  The report's Python reference uses this variant.  For a first object-theory implementation, the dual-tree form has the advantage that the logical core is just equality propagation on a finite graph.

\section{Object-theory dependency graph}
\label{app:dependencies}

\begin{center}
\begin{tabularx}{\textwidth}{@{}lX@{}}
\toprule
Stage & Depends on \\
\midrule
Point determinant identities & \code{Core.Integers}, \code{Core.LinearAlgebra}, \code{Core.Tactic} \\
Path folds and reversal & \code{Core.Definitions}, \code{Core.Pairing}, quantifier-free induction \\
Path shoelace lemmas & point identities, \code{Core.BoundedSums} \\
Disk boundary cancellation & path folds, triangle incidence certificate \\
Conway period-area theorem & path centrality, translated-pair lemma, disk cancellation \\
Lattice coordinates and clipping & rational arithmetic, \code{Core.Order}, \code{Core.Search} \\
Torus refinement & clipping, finite sorting, pairing/list infrastructure \\
Constant multiplicity & edge incidence, dual spanning tree, induction \\
Multiplicity one & period-area theorem, bounded area sums, positive cancellation \\
Final criterion & all previous stages, propositional composition and universal closure \\
Native export & final \(\Prf\,\PRA\) theorem, \code{pra\_iff\_native} \\
\bottomrule
\end{tabularx}
\end{center}

The dependency graph intentionally keeps topology downstream of exact arithmetic.  This makes failures local: the affine and area theorem can be promoted to the checked derivation library before the full torus constructor is finished.

\section{Theorem and artifact ledger}
\label{app:ledger}

\begin{longtable}{@{}p{0.12\textwidth}p{0.53\textwidth}p{0.25\textwidth}@{}}
\toprule
ID & Claim/artifact & Status \\
\midrule
\endhead
C1 & Composition of half-turns is translation & Proved in \cref{lem:halfturnalgebra}. \\
C2 & Four turned-orbit path identities & Proved in \cref{lem:orbitpairs}. \\
C3 & Open-path shoelace laws & Proved in \cref{lem:Jlaws}. \\
C4 & Central path endpoint determinant & Proved in \cref{lem:Jcentral}. \\
C5 & \(\det(t,u)=2\Area(T)>0\) & Proved in \cref{thm:periodarea}. \\
C6 & Projected two-tile boundary is zero & Proved in \cref{lem:boundaryzero}. \\
C7 & Closed nonnegative top chain has constant multiplicity & Proved in \cref{lem:constantmultiplicity}. \\
C8 & Multiplicity is one & Proved by \cref{eq:kcancel}. \\
C9 & Explicit family tiles the plane & Proved in \cref{thm:mainPL}. \\
P1 & \code{code/conway\_core.py} & Exact rational reference implementation. \\
P2 & \code{code/test\_conway\_core.py} & Seven passing tests at delivery. \\
P3 & Worked torus certificate & \(58\) pieces, exact area \(1\), zero positive-area overlap. \\
L1 & \code{lean/ConwayCriterionBlueprint.lean} & Uncompiled definitions/targets; no admitted facts. \\
L2 & Empty-context \(\Prf\,\PRA\) theorem & Planned, not claimed complete. \\
\bottomrule
\end{longtable}

\section{Reproduction commands}
\label{app:reproduce}

From the package root:

\begin{lstlisting}
(cd code && python3 -m unittest -v test_conway_core.py)
python3 code/conway_core.py examples/zigzag_conway_tile.json \
  --output examples/zigzag_conway_tile_report.json
python3 figures/make_figures.py
bash build.sh
\end{lstlisting}

The exact kernel depends only on the Python standard library.  Figure generation requires Matplotlib.  The PDF build requires a conventional TeX Live installation with the packages in the preamble.
